\section*{結言}
本日のX上では、AI事業の売上や将来予測、AI生成研究への批判など、注目度の高い話題ほど数字や引用の前提が揺れていた。企業公式の新モデル発表が少なかったこともあり、一次情報に遡れない投稿を断定的なニュースに変えないことが重要となる。

一方で、文書処理や音声検出の製品化、推論資源の配分、エージェントの実行環境・評価・決済という話題は、モデル単体の性能から運用システム全体へ関心が移る流れを示した。Haiku 5.5やSalukiの比較も、価格表やベンチマークの数字だけでなく、実際のトークン消費、量子化、タスク適合性を確かめる必要性を浮かび上がらせている。未確認事項と観測窓外の発表は本号でも明示した。
